% Continuation of the existing direct-product and bounded-power treatment.
\section{Reduced products and further Boolean-power constructions}
\label{sec:p03-fv-reduced-composition}

The Boolean reduction of Definition~\ref{def:p03-fv-reduction} also
computes truth in a reduced product. Its Boolean algebra is now a quotient
of the algebra of index sets. A partition in that quotient must be realized
by actual index sets before it can supply a product witness. This finite
lifting step is the difference between a complete construction and an
unjustified appeal to the ultraproduct theorem. The classical construction
below is separate from the real-valued metric construction of
\cite[\S2.3 and \S3]{Ghasemi2016MetricFVv3}; the latter explicitly records
the classical quotient-Boolean-algebra theorem on p.~1. The notation
\(\Omega,I,A_\gamma,\pi_I\) retains that source's index/ideal roles.
In particular \(I\) here is an ideal, not the index set denoted \(I\)
in Definition~\ref{def:p03-fv-product}.

\begin{hbdefinition}{def:p03-fv-reduced-standing}{Classical reduced products over an ideal}
Let \(\Omega\) be a nonempty set and let \(I\) be a proper ideal of
\(\mathcal P(\Omega)\): it contains \(\varnothing\), is closed under
subsets and finite unions, and does not contain \(\Omega\).
The dual proper filter is
\(F_I=\{S\subseteq\Omega:\Omega\setminus S\in I\}\).
All \(A_\gamma\), \(\gamma\in\Omega\), have the same finitary
many-sorted first-order signature \(L\), with nonempty sort domains.
The construction of the full sorts and the witness selection below take
place in \(\mathsf{ZF}+\mathsf{AC}\).

In each sort identify product elements \(a,b\) when
\(\{\gamma:a(\gamma)\ne b(\gamma)\}\in I\). Write \(\pi_I(a)\)
for the class and \(A_I=\prod_I A_\gamma\) for the quotient structure.
Constants and functions descend from coordinatewise interpretation.
For every relation symbol, including a nullary one, set
\[
 A_I\models R(\pi_I(\bar a))
 \quad\Longleftrightarrow\quad
 \{\gamma:A_\gamma\not\models R(\bar a(\gamma))\}\in I.
\]
This is the ordinary two-valued interpretation. For discrete metric
structures with distance zero on equality and one on inequality, and a
predicate valued zero on truth and one on falsity, the source's
\(I\)-limsup atomic interpretation has exactly this zero condition.
It does not identify arbitrary non-atomic classical satisfaction with a
coordinate limsup.

In the Boolean algebra \(\mathcal P(\Omega)/I\), two subsets have the
same class when their symmetric difference belongs to \(I\).
Write \([S]_I\) for their class. Boolean operations are induced by the
ordinary finite set operations. For a factor formula define
\[
 X_\psi(\bar a)=\{\gamma:A_\gamma\models\psi(\bar a(\gamma))\},
 \qquad [X_\psi(\bar a)]_I\in\mathcal P(\Omega)/I.
\]
The local value refers to satisfaction in the factors, not in \(A_I\).
\end{hbdefinition}

\begin{hblemma}{thm:p03-fv-finite-partition-lift}{proved=annex:p03-fv-finite-partition-lift}{Lifting a finite Boolean partition}
Let \(h:B\to C\) be a surjective Boolean-algebra homomorphism and
let \(c_1,\ldots,c_n\) be a finite partition of \(1_C\), with
\(n\ge1\) and zero cells allowed. There is a partition
\(b_1,\ldots,b_n\) of \(1_B\) such that \(h(b_k)=c_k\).
Consequently every finite partition in \(\mathcal P(\Omega)/I\)
has representatives that partition \(\Omega\) exactly, not merely
modulo \(I\).
\end{hblemma}

\begin{hblemma}{thm:p03-fv-reduced-well-defined}{proved=annex:p03-fv-reduced-well-defined}{Quotient interpretations and local formula classes}
The element equivalence, the coordinatewise functions and the declared
relations define an \(L\)-structure. Every local formula class
\([X_\psi(\bar a)]_I\) is independent of representatives of the
finite parameter tuple. Terms evaluate as their coordinatewise classes.
For every atomic \(\alpha\),
\[
 A_I\models\alpha(\pi_I(\bar a))
 \quad\Longleftrightarrow\quad
 [X_\alpha(\bar a)]_I=1.
\]
This does not assert the same equivalence for arbitrary \(\alpha\).
\end{hblemma}

\begin{hbdefinition}{def:p03-fv-reduced-reduction}{The quotient-index Boolean condition}
For an input \(\varphi(\bar x)\), use the finite factor list and
Boolean formula \((\psi_1,\ldots,\psi_m;\Theta_\varphi)\) computed
in Definition~\ref{def:p03-fv-reduction}. Its atomic clause is \(X=1\);
falsity is \(0\ne0\); negation negates the Boolean \emph{formula};
conjunction concatenates the lists after disjoint renaming. Boolean
quantifiers now range over \(\mathcal P(\Omega)/I\).
In the existential clause the local formulas are the complete signed
patterns \(\delta_\sigma(\bar x,y)\) and their availability formulas
\(\eta_\sigma(\bar x)=\exists y\,\delta_\sigma(\bar x,y)\),
for \(\sigma\in\{0,1\}^m\). The condition is
\[
 \exists(Z_\sigma)_\sigma\left[
 \operatorname{Part}_{\{0,1\}^m}((Z_\sigma)_\sigma)
 \land\bigwedge_\sigma Z_\sigma\le Y_\sigma
 \land\Theta\left(\bigvee_{\sigma_1=1}Z_\sigma,\ldots,
                         \bigvee_{\sigma_m=1}Z_\sigma\right)\right],
\]
with \(Y_\sigma=[X_{\eta_\sigma}(\bar a)]_I\).
The partition cells and the availability values are Boolean classes.
They must not be confused with chosen subsets or factor witnesses.
All Boolean binders and factor binders are freshened capture-avoidantly
as in the direct-product construction. The one empty pattern when
\(m=0\) is retained.
\end{hbdefinition}

\begin{hbtheorem}{thm:p03-fv-reduced-reduction}{proved=annex:p03-fv-reduced-reduction}{Composition theorem for arbitrary proper-ideal reduced products}
For every fixed first-order formula the effective finite reduction just
specified satisfies
\[
 A_I\models\varphi(\pi_I(\bar a))
 \quad\Longleftrightarrow\quad
 \mathcal P(\Omega)/I\models
 \Theta_\varphi([X_{\psi_1}(\bar a)]_I,\ldots,
                           [X_{\psi_m}(\bar a)]_I).
\]
Its existential implication constructs an actual representative witness
outside one \(I\)-small exceptional set, fills that set with arbitrary
elements, and then passes to its class. No ultrafilter alternative, Boolean
completeness, atomlessness, or countability of \(\Omega\) is assumed.
The case \(I=\{\varnothing\}\) is the direct-product theorem.
An improper ideal and an empty index set are outside this proper-filter
formulation; the separate empty-direct-product convention remains as stated
in Definition~\ref{def:p03-fv-standing}.
\end{hbtheorem}

\begin{hbcorollary}{thm:p03-fv-reduced-elementarity}{proved=annex:p03-fv-reduced-elementarity}{Factorwise elementarity and theory transfer}
For the same \(\Omega,I\), if \(A_\gamma\preccurlyeq B_\gamma\)
for every index, the induced inclusion of reduced products is elementary.
If instead \(A_\gamma\equiv B_\gamma\) for every index, their reduced
products are elementarily equivalent. In the inclusion assertion each sort
is included as a subset; the corresponding statement for specified elementary
embeddings uses their coordinatewise map. The Boolean coefficient algebra
is held fixed. Merely changing to an isomorphic or elementarily equivalent
coefficient algebra does not provide a map between arbitrary varying-factor
products without additional data.
\end{hbcorollary}

\begin{hbtheorem}{thm:p03-fv-reduced-example}{proved=annex:p03-fv-reduced-example}{A worked witness and a failure of unrestricted \L o\'s transfer}
For binary \(R\) and unary \(P\), let
\[
 U=[X_{\exists y(R(x,y)\land\neg P(y))}(a)]_I,
 \qquad V=[X_{\exists y(R(x,y)\land P(y))}(a)]_I.
\]
Then \(A_I\models\exists y(R(x,y)\land\neg P(y))(\pi_I(a))\)
exactly when \(U\vee V=1\) and \(U\ne0\).
For example, on \(\Omega=\mathbb N\) modulo the finite-set ideal,
let every factor have domain \(\{0,1\}\), let \(R\) hold on every pair,
and let \(P\) be empty at even indices and universal at odd indices.
Every tuple in the reduced product fails \(P\); its negation is true,
although the local negation truth set is the even numbers and is not
cofinite. Thus the arbitrary-formula assertion
\(A_I\models\psi\Longleftrightarrow X_\psi\in F_I\) is false
for this proper filter. The atomic criterion and composition theorem
remain true. An ultrafilter excludes this particular obstruction because
its quotient Boolean algebra has exactly two elements.
\end{hbtheorem}

\subsection{Product, iteration, and quotient laws for bounded Boolean powers}

\begin{hbconvention}{def:p03-fv-source-power-standing}{The bounded algebra power and its Stone notation}
The following laws concern the bounded power \(A[B]^*\) of
\cite[IV\S5, Theorem~5.4 and Exercises~1--3, pp.~141--145]{BurrisSankappanavar2012},
not arbitrary full function powers. In that source \(B^*\) is the compact
Stone space of ultrafilters of \(B\), and \(A[B]^*\) consists of
continuous maps \(B^*\to A\), with \(A\) discrete. Here \(A\)
is a nonempty algebra of a finitary functional signature, just as in the
source. Stone representation uses the Boolean Prime Ideal Theorem.
The equivalent partition presentation and its finite refinements were
proved in Theorem~\ref{thm:p03-fv-power-representations}.
Write \(C(X,A)\) for the algebra of continuous functions from a compact
Stone space \(X\) to the discrete algebra \(A\), with pointwise operations.
It is the same bounded construction when \(X=B^*\).
The basic clopen attached to \(b\in B\) is
\(N_b=\{U\in B^*:b\in U\}\), in the source's notation.
\end{hbconvention}

\begin{hbtheorem}{thm:p03-fv-source-power-products}{proved=annex:p03-fv-source-power-products}{The bounded-power product laws}
If \(B\) is nontrivial, evaluation at every point of \(B^*\) is
surjective, and constants embed \(A\) in \(A[B]^*\). Thus it is
a subdirect power of \(A\). Also
\[
 A[\mathbf 2]^*\cong A,
 \qquad
 A[B_1\times B_2]^*\cong A[B_1]^*\times A[B_2]^*,
 \qquad
 (A_1\times A_2)[B]^*\cong A_1[B]^*\times A_2[B]^*.
\]
The coefficient-product isomorphism restricts a function to the two
clopen components of \((B_1\times B_2)^*\); the base-product
isomorphism takes its two coordinate functions. The trivial coefficient
algebra \(\mathbf1\) has empty Stone space and gives a one-element
power. In that case the evaluation/subdirect assertion has no nonempty
index set and constants need not embed a nontrivial base.
\end{hbtheorem}

\begin{hbdefinition}{def:p03-fv-stone-product-coefficients}{The source's coefficient operation \(B_1*B_2\)}
For a Boolean space \(X\), write \(X^*\) for its Boolean algebra of
clopen subsets; for a Boolean algebra \(B\), \(B^*\) is its Stone
space. These are opposite typed uses of the star, not exponentiation.
The source's operation is
\[
                 B_1*B_2=(B_1^*\times B_2^*)^*.
\]
This is distinct from the direct product \(B_1\times B_2\), whose
Stone space is a disjoint union rather than a Cartesian product.
The inner algebra \(A[B_1]^*\) is given the discrete topology when
it is used as the base for the next bounded power.
\end{hbdefinition}

\begin{hbtheorem}{thm:p03-fv-iterated-power}{proved=annex:p03-fv-iterated-power}{Iteration of bounded powers}
There are natural algebra isomorphisms
\[
       (A[B_1]^*)[B_2]^*\cong A[B_1*B_2]^*
              \cong (A[B_2]^*)[B_1]^*.
\]
The first map sends \(g:B_2^*\to A[B_1]^*\) to the function
\((x,y)\mapsto g(y)(x)\). Conversely a continuous function on the
product has only finitely many distinct continuous slices, and their
slice classes are clopen in \(B_2^*\). The inverse therefore really
belongs to the \emph{bounded} iterated power; an arbitrary function-space
currying identity would not establish this condition.
\end{hbtheorem}

\begin{hbdefinition}{def:p03-fv-power-filter-quotient}{The source's filter congruence}
Let \(F\) be a filter of a Boolean algebra \(B\), with the improper
filter \(F=B\) permitted in this source presentation.
Write \(q_F:B\to B/F\) for the quotient by
\[
 b\equiv_F c\quad\Longleftrightarrow\quad
              (b\wedge c)\vee(\neg b\wedge\neg c)\in F.
\]
The closed subset \(F^*\subseteq B^*\) consists of the ultrafilters
containing \(F\). For power elements \(a,b\), the source defines
\[
             (a,b)\in\theta_F
             \quad\Longleftrightarrow\quad
             [\![a=b]\!]\supseteq F^*,
\]
where the bracket now denotes their clopen equalizer in \(B^*\).
Equivalently its corresponding Boolean element belongs to \(F\).
The finite-partition presentation uses the same equalizer as the join
of all common-refinement cells whose labels agree.
\end{hbdefinition}

\begin{hbtheorem}{thm:p03-fv-power-quotient}{proved=annex:p03-fv-power-quotient}{Quotient of a bounded power}
The relation \(\theta_F\) is an algebra congruence and
\[
                    A[B]^*/\theta_F\cong A[B/F]^*.
\]
On finite labeled partitions the inducing surjection sends each coefficient
\(b_r\) to \(q_F(b_r)\) without changing its base label. Its kernel
is exactly \(\theta_F\), and finite partition lifting proves
surjectivity. Hence this is a proved algebra quotient, not a mere
analogy with a reduced product. No claim of elementary preservation under
this arbitrary quotient is made. If \(F=B\), its closed Stone subspace
is empty, the congruence is universal, and both quotient powers have one
element, including when \(B\) itself is trivial.
\end{hbtheorem}

\begin{hbtheorem}{thm:p03-fv-power-cardinality}{proved=annex:p03-fv-power-cardinality}{Size of an infinite bounded Boolean power}
Assume Choice, \(A\) has at least two elements and \(B\) is a
nontrivial Boolean algebra. If either cardinal \(|A|\) or \(|B|\)
is infinite, then
\[
                 |A[B]^*|=|A|\cdot|B|=\max(|A|,|B|).
\]
The restriction to the bounded finite-partition power matters: the full
Cartesian power can be much larger. The assumption on \(A\) cannot
be removed, since the power of a one-element algebra has one element
even over an infinite \(B\).
\end{hbtheorem}
